% !PRINTED 593
\weberpage{593}
%<<CONT
\weq{1}{\mu = \overset{\varrho}{\Pi}\,\boldsymbol{s}\!\left(\frac{2\varrho}{\mu}\right)}
kommt eine Zahl $2\varrho$ vor, die mit $\mu\eta$ nach dem Modul $\mu$ kongruent
 ist. Folglich ist $\boldsymbol{s}(\eta)$ ein Teiler von $\mu$ und mithin relativ
prim zu $\mathfrak{m}$, wie bewiesen werden sollte.

Nehmen wir jetzt an, da\ss{} in der Formel (11), \S\,158
\weq{2}{\tau_0 = \overset{\xi}{\Pi}\,\boldsymbol{s}\!\left(\frac{2\alpha\xi}{\mu}\right)}
$\mathfrak{m}$ ein Primideal des K\"orpers $\Omega$ sei und setzen demgem\"a\ss{}
$\mathfrak{m} = \mathfrak{p}$, so da\ss{} $\xi$ ein volles Restsystem nach dem Modul $\mathfrak{p}$ mit
Ausschlu\ss{} der Null durchl\"auft.

Nehmen wir dann $\mu$ so an, da\ss{} es nur durch die erste Potenz
von $\mathfrak{p}$ teilbar und durch ein beliebig gew\"ahltes anderes Primideal
$\mathfrak{p}'$ nicht teilbar ist, so kann $\tau_0$ nach dem Satz 6.\ nicht durch $\mathfrak{p}'$
teilbar sein. Nun ist $\tau_0$ nur von dem Ideal $\mathfrak{p}$ abh\"angig und
unabh\"angig davon, wie im \"ubrigen die Zahl $\mu$ genommen ist.
Folglich ist $\tau_0$ nach dem Satz 6.\ durch kein von $\mathfrak{p}$ verschiedenes
Primideal teilbar. Es kann aber auch $\mathfrak{p}$ nicht in einer h\"oheren
als der ersten Potenz in $\tau_0$ aufgehen, denn die Faktoren des
Produktes (1) kommen alle unter den Faktoren des Produktes (2)
vor, und folglich ist $\mu$ durch $\tau_0$ teilbar. Demnach k\"onnen wir
geradezu
\weq{3}{\tau_0 = \mathfrak{p}}
setzen, und da $\tau_0$ eine Zahl im K\"orper $\mathfrak{L}$ ist, so folgt:

7. \gesperrt{Jedes ungerade Primideal des K\"orpers $\Omega$ ist ein
K\"orper $\mathfrak{L}$, ein Hauptideal.}
% [sic] printed ``ist ein Körper $\mathfrak L$, ein Hauptideal''; sense requires ``ist im Körper $\mathfrak L$ ein Hauptideal''.

Auch in der Annahme, da\ss{} $\mu$ relativ prim zu $\mathfrak{a}$ sein sollte,
liegt keine Einschr\"ankung dieses Satzes. Denn wir k\"onnen bei
gegebenem $\mathfrak{p}$ unter den \"aquivalenten Formen $(a, b, c)$, ohne $\lambda$ zu
\"andern, eine ausw\"ahlen, bei dem diese Forderung erf\"ullt ist.

Noch nicht bewiesen ist aber hierdurch, da\ss{} $\mathfrak{p}$ auch im
Klassenk\"orper $\mathfrak{K}(\Delta)$ selbst ein Hauptideal ist. Wenigstens folgt
dies nur in dem Fall $\Delta \equiv 1\ (\mathrm{mod}\ 8)$, wo $\mathfrak{K}$ mit $\mathfrak{L}$ identisch
ist\footnote{Vgl.\ hierzu \S\,122, wo auch der Fall der Primzahl 2 erledigt ist.}.

Nach dem Satz 3., \S\,157 sind die Faktoren des Produktes (2)
alle miteinander assoziiert. Die Anzahl dieser Faktoren ist
$N(\mathfrak{p}) - 1$, und demnach k\"onnen wir auch setzen:
\weq{4}{\mathfrak{p} = \Bigl[\boldsymbol{s}\!\left(\frac{2\alpha}{\mu}\right)\Bigr]^{N(\mathfrak{p})-1}.}
